
\chapter{Opis projektnog zadatka}
		
		Cilj ovog projekta je izgradnja web aplikacije koja će omogućiti ljudima/pacijentima lakše prijavljivanje na slobodne termine za medicinsku rehabilitaciju i fizikalnu terapiju te praćenje njihovog zdravstvenog napretka. Također će omogućiti zaposlenicima ustanove da otkazuju ili mijenjaju termine ovisno o raspoloživosti prostorija, opreme i osoblja pritom imajući mogućnost uvida u pacjentove prošle terapije i napredak. Vrijeme provođenja rehabilitacije je svakim radnim danom od ponedjeljka do petka od 8:00 do 20:00 sati.
		
		\noindent Aplikacija će razlikovati tri vrste korisnika: 
		\begin{packed_item}
			
			\item  pacijenta
			\item  liječnika - djelatnika zdravstvene ustanove
			\item  administratora - djelatnika zdravstvene ustanove
		\end{packed_item}
		
		Prilikom pokretanja web aplikacije svaki korisnik unosi svoju e-mail adresu i lozinku. Ovisno o vrsti korisnika biti će preusmjereni na različite stranice. Svaki korisnik imati će mogućnost mijenjanja nekih osobnih podataka i lozinke. Kao i promjene lozinke u slučaju da je zaboravljena.
		
		\textit{Pacijent} se samostalno prijavljuje u sustav. U slučaju da još ne postoji imati će opciju registracije. Za registraciju mora unijeti: 
		\begin{packed_item}
			\item ime
			\item prezime
			\item e-mail adresu
			\item MBO - Matični Broj Osiguranika
			\item broj telefona
			\item lozinku
		\end{packed_item}
		Prilikom registracije pomoću MBO-a provjerava se ako korisnik postoji u središnjem informacijskom sustavu zdravstvene zaštite.
		Nakon prijave korisnik je preusmjeren na početnu stranicu gdje može prikazati svoje termine, terapije i prijavljivati iste. U terminima se prikazuju protekli i budući termini, za protekle termine piše ako je pacijent došao te komentari djelatnika ustanove o napretku. Termin se dobiva nakon što se odobri zahtjev za jednim. No prije nego pacijent može birati termin mora kreirati terapiju, u slučaju da nema ni jednu aktivnu. Terapiju kreira unoseći podatke:
		\begin{packed_item}
			\item ime i prezime liječnika koji ga je uputio na terapiju
			\item tip rehabilitacije
			\item opis oboljenja 
			\item opis postupka liječenja
		\end{packed_item}
		Nakon što preda podatke o terapiji, provjerava se status liječnika u imeniku liječnika. U slučaju da uneseni liječnik ne postoji, zbog krivo unesenih podataka sustav će upozoriti pacijenta o tome i imati će priliku ispraviti podatke. Nakon uspješno kreirane terapije pacijent može kreirati termin. 
		Termin se sastoji od:
		\begin{packed_item}
			\item željenog vremena 
			\item vrste/tipa rehabilitacije
			\item reference na terapiju
			\item komentara za liječnika
		\end{packed_item}
		Pacijent, nakon što je kreirao terapiju, unutar nje kreira termin odabirući pritom prostoriju i željeno vrijeme. Ukoliko se željeni termin nije popunio, termin će automatski biti odobren. Pacijentu će o odobrenom terminu, te o svim mogućim promjenama biti obaviješten e-mailom.
		
		\textit{Liječnik} nakon prijavljivanja u sustav ima pregled svih pacijenata i njihovih podataka. Imati će opciju pretraživanja pacijenta. Klikom na prikaz bit će mu prikazani svi termini odabranog pacijenta i detalji o njima, tu će imati opciju evidentirati dolazak pacijenta i zabilježiti komentare vezane uz napredak. Liječnik koji je evidentirao pacijenta dodati će se u sustav kao liječnik koji je vodio taj termin terapije. Također svaki djelatnik ustanove imati će mogućnost promjene termina, o čemu će pacijent biti obaviješten e-mailom. 
		
		\textit{Administrator} ima pregled svih pacijenata i djelatnika. Uz ovlasti koje imaju liječnici, administrator pri zaposlenju novog liječnika izrađuje korisnički račun za njega s inicijalnom lozinkom. Također nakon prestanka radnog odnosa administrator može ukloniti tog liječnika. Administrator definira sve što je potrebno za ispravan rad sustava, dakle može mijenjati dostupnost opreme i unositi novu opremu, mijenjati dostupnost prostorija i unositi nove. Također ako dođe do promjena u dostupnosti prostorija ili opreme.
		
		---
		
		Ovim projektom smanjio bi se opseg posla djelatnika ustanove i olakšao proces prijave na rehabilitaciju pacijentima, umjesto prijavljivanja u živo i rješavanja papirologije pacijentima će biti omogućeno automatsko prijavljivanje na slobodne termine, a liječnicima evidentiranje istih. 
		
		---
		
		Odgovarajući primjer slične aplikacije na istu temu nije pronađen, moguće, zbog zatvorenosti ovakvog tipa aplikacije prema javnosti, ali primjer sa sličnim funkcijama iako ne u istu svrhu je booking.com (\ref{fig:booking}). Na početnoj stranici \textit{booking}-a nailazimo na opcije prijave i registracije, iako za početak korištenja same stranice to nije nužno kao kod ove aplikacije. Korisnik \textit{booking}-a rezervira apartman, dok korisnik aplikacije za rehabilitaciju rezervira svoj termin za terapiju. Kao što korisnik na \textit{booking}-u unosi grad u kojem želi rezervirati apartman, korisnik u ovoj aplikaciji odabire tip rehabilitacije za koju se želi prijaviti. U oba slučaja kod prijave na željeni termin ili rezervacije apartmana korisnik dobiva neke ponuđene termine/opcije ili može samostalno odabrati vrijeme željenog termina/rezervacije. 
		
		\begin{figure}[H]
			\includegraphics[scale=0.4]{slike/slicna_aplikacija.PNG} %veličina slike u odnosu na originalnu datoteku i pozicija slike
			\centering
			\caption{Primjer slične aplikacije}
			\label{fig:booking}
		\end{figure}
		
		---
		
		Aplikacija će se moći proširiti na način da su svi termini, prostorije, oprema i djelatnici vidljivi su administratoru koji im može mijenjati status za određeno vrijeme, na primjer kada liječnik ode na godišnji odmor može promijeniti njegov status u neaktivan za to razdoblje ili kada cijela ustanova ima sastanak onemogućiti sve termine za vrijeme tog sastanka. Također može se proširiti uvođenjem statistike koja bi ustanovi omogućila bolje korištenje resursa i nabavku novih. Primjerice ako je veća potražnja za kardiovaskularnim terapijama da povećaju broj opreme i liječnika specijaliziranih za taj tip terapije.
		
	
		\eject
		
				
	